% =====================================================================
%  macros.tex: notation
% =====================================================================

% ---------- number systems and scalar field ----------
\newcommand{\N}{\mathbb{N}}
\newcommand{\Z}{\mathbb{Z}}
\newcommand{\Q}{\mathbb{Q}}
\newcommand{\R}{\mathbb{R}}
\newcommand{\C}{\mathbb{C}}
\newcommand{\K}{\mathbb{K}}

% ---------- delimiters ----------
\DeclarePairedDelimiter{\abs}{\lvert}{\rvert}
\DeclarePairedDelimiter{\norm}{\lVert}{\rVert}
\DeclarePairedDelimiter{\set}{\{}{\}}

% ---------- spaces ----------
\newcommand{\Cont}[1]{C[#1]}
\newcommand{\Contk}[2]{C^{#1}[#2]}
\newcommand{\lp}[1]{\ell_{#1}}
\newcommand{\Lp}[1]{L_{#1}}

% ---------- metric and topological notation ----------
\newcommand{\cl}[1]{\overline{#1}}

% ---------- miscellaneous ----------
\newcommand{\eps}{\varepsilon}
\newcommand{\dd}{\mathrm{d}}

% ---------- bounded operators (added with the lecture of 2026-09-24) ----------
% Script B keeps operator spaces apart from balls B(x, eps).
\newcommand{\Bop}[2]{\mathcal{B}(#1, #2)}
% Shorthand for operators of a space into itself (lecture of 2026-10-01).
\newcommand{\Bopone}[1]{\mathcal{B}(#1)}

% ---------- inner product spaces (added with the lecture of 2026-10-08) ----------
% First slot linear, second slot conjugate linear.
\DeclarePairedDelimiterX{\ip}[2]{\langle}{\rangle}{#1, #2}
\newcommand{\orth}[1]{#1^{\perp}}
